% GDZ PPN236018647,native585lower–589upper;HKUST586–590.
\NumberedItem{116}
\noindent\hspace*{\parindent}Impossibile esse, ut sectio circuli ad aequationes inferiores, quam theoria nostra suggerit, reducatur, demonstratum.
\par\noindent\hfill Brunov[ici, 1801] Apr. 6.
\par\medskip\begingroup\small
\noindent\hspace*{\parindent}Dieselbe Aussage findet sich auch in den \textit{Disquisitiones arithmeticae} (1801) und zwar im art. 365. (Werke I, S. 462) für den Fall, wo die Anzahl \(n\) der Teile, in die der Kreis zu teilen ist, eine Primzahl, im art. 366. (ebenda, S. 463) für die Fälle, wo diese Anzahl eine Primzahlpotenz und eine beliebige zusammengesetzte Zahl ist. Wahrscheinlich sind diese Bemerkungen ebenso wie die entsprechenden, ihnen im
\par\endgroup

\clearpage\SourceLeaf{586}\ArithmeticHead{557}{R}{1800 OCT. 25—1801 APR. 6.}
\par\begingroup\small
\renewcommand{\thefootnote}{*)}
\noindent art. 336. (ebenda, S. 413) voraufgehenden erst während der Drucklegung eingefügt worden\footnote{Vergl. auch den oben S. 121 abgedruckten Brief von \textsc{Gauss} an \textsc{Gerling} vom 6. Januar 1819, wo es (oben S. 125) heißt: »Alles hängt dabei von den Factoren der Zahl \(\frac12(p-1)\) ab; ist diese Zahl eine Potenz von 2, z. B. \(p=3,\ 5,\ 17,\ 257,\ 65537\), so kommen bloß quadratische Gleichungen vor; hingegen z. B. für \(p=31\), wo \(\frac12(p-1)=3\mathord{.}5\), ist eine cubische und eine Gleichung vom 5. Grade unausweichlich«.}. Ein Beweis für \textsc{Gauss}' Behauptung ist aber weder in seinen Veröffentlichungen noch in seinem Nachlasse enthalten, auch ist ein Beweis, der nur solche Hilfsmittel benutzt, wie \textsc{Gauss} sie im Jahre 1801 sonst angewandt hat, bisher nicht veröffentlicht worden. Es soll darum hier ein elementarer Beweis ohne \textsc{Galois}sche Theorie gegeben werden; er gründet sich auf den folgenden \textls[120]{Satz}:
\par\noindent\hspace*{\parindent}Gegeben sei eine Gleichung \(J(x)=0\) vom \(n\)-ten Grade, deren Koeffizienten dem Rationalitätsbereiche \(P\) angehören und die in \(P\) irreduzibel ist. Eine Wurzel \(\xi_1\) dieser Gleichung sei rational mit Koeffizienten aus \(P\) dargestellt durch die Größen \(\rho_1,\rho_2,\ldots,\rho_k\), wo \(\rho_1\) einer Gleichung \(X_1(x)=0\) genügt, deren Koeffizienten \(P\) angehören, und die in \(P\) irreduzibel ist, \(\rho_2\) einer Gleichung \(X_2(x)=0\), deren Koeffizienten dem durch Adjunktion von \(\rho_1\) erweiterten Rationalitätsbereiche \((P,\rho_1)\) angehören und die in diesem Bereiche irreduzibel ist u.s.w., endlich \(\rho_k\) einer Gleichung \(X_k(x)=0\), deren Koeffizienten dem Rationalitätsbereiche \((P,\rho_1,\rho_2,\ldots,\rho_{k-1})\) angehören und die in diesem Bereiche irreduzibel ist. \textls[120]{Dann ist das Produkt der Grade aller Gleichungen \(X_1(x)=0,\ X_2(x)=0,\ldots,\ X_k(x)=0\) durch den Grad \(n\) von \(J(x)=0\) teilbar.}
\par\noindent\hspace*{\parindent}Adjungiert man nämlich dem Bereiche \(P\) der Reihe nach \(\rho_1,\rho_2,\ldots\), so sei \(\rho_\alpha\) die erste dieser Größen, durch deren Adjunktion \(J(x)=0\) reduzibel wird; die Wurzel \(\xi_1\) genügt dann in dem Rationalitätsbereiche \((P,\rho_1,\rho_2,\ldots,\rho_\alpha)\) einer irreduziblen Gleichung, deren Grad \(n'\) kleiner als \(n\) ist. Es sei \(h_\alpha\) der Grad der in \((P,\rho_1,\rho_2,\ldots,\rho_{\alpha-1})\) irreduziblen Gleichung \(X_\alpha(x)=0\), der \(\rho_\alpha\) genügt; dann ist nach dem Satze von \textsc{Kronecker-Kneser}\renewcommand{\thefootnote}{**)}\footnote{Dieser Satz lautet wie folgt: Sind \(J_1(x)=0\) und \(J_2(x)=0\) zwei im Rationalitätsbereiche \(P\) irreduzible Gleichungen der Grade \(n_1\) und \(n_2\) mit den Wurzeln \(\eta_1\) und \(\eta_2\) und sind die Grade der irreduziblen Gleichungen, denen \(\eta_1\) nach Adjunktion von \(\eta_2\) und \(\eta_2\) nach Adjunktion von \(\eta_1\) genügen, \(n'_1\) und \(n'_2\), so ist \(\frac{n_1}{n'_1}=\frac{n_2}{n'_2}\). Siehe \textsc{A. Kneser}, Mathem. Annalen \textbf{30}, 1887, S. 195, \textsc{Crelles} Journal für Mathem. \textbf{106}, 1890, S. 51, \textsc{O. Hölder}, Mathem. Annalen \textbf{38}, 1891, S. 309, \textsc{G. Frobenius}, ebenda \textbf{70}, 1911, S. 457. Der Beweis, den \textsc{Kneser} im 106. Bande des \textsc{Crelles}chen Journals für diesen Satz gibt, erfordert nur die einfachsten Hilfsmittel. Vgl. auch \textsc{A. Loewy}, Jahresbericht der deutschen Mathematikervereinigung 1917.} \(\frac n{n'}=\frac{h_\alpha}{h'_\alpha}\), wenn \(h'_\alpha\) der Grad derjenigen irreduziblen Gleichung ist, die nach Adjunktion von \(\xi_1\) zum Rationalitätsbereiche \((P,\rho_1,\rho_2,\ldots,\rho_{\alpha-1})\) durch \(\rho_\alpha\) befriedigt wird. Bei weiterer Adjunktion der auf \(\rho_\alpha\) folgenden Größen bleibe die Gleichung \(n'\)-ten Grades für \(\xi_1\) irreduzibel bis zur Adjunktion von \(\rho_\beta\). In dem Bereiche \((P,\rho_1,\ldots,\rho_\alpha,\ldots,\rho_\beta)\) genüge \(\xi_1\) einer irreduziblen Gleichung vom Grade \(n''<n'\). Bezeichnet dann \(h_\beta\) den Grad der im Bereiche \((P,\rho_1,\ldots,\rho_\alpha,\ldots,\rho_{\beta-1})\) irreduziblen Gleichung \(X_\beta(x)=0\), der \(\rho_\beta\) genügt, und \(h'_\beta\) den Grad der irreduziblen Gleichung, die sich für \(\rho_\beta\) im Bereiche \((P,\rho_1,\ldots,\rho_\alpha,\ldots,\rho_{\beta-1},\xi_1)\) ergibt, so ist \(\frac{n'}{n''}=\frac{h_\beta}{h'_\beta}\). Derart fortfahrend erhält man die weiteren Beziehungen
\[
\frac{n''}{n'''}=\frac{h_\gamma}{h'_\gamma},\quad\frac{n'''}{n^{(4)}}=\frac{h_\delta}{h'_\delta},\quad\cdots.
\]
\noindent Da \(\xi_1\) schließlich rational in \(\rho_1,\rho_2,\ldots,\rho_k\) wird, also einer Gleichung ersten Grades genügt, so ergibt sich die Beziehung \(\frac{n^{(t-1)}}1=\frac{h_\tau}{h'_\tau}\), wo \(t\leqq k\) und \(h_\tau\) der Grad jener irreduziblen Gleichung \(X_\tau(x)=0\) ist, deren Wurzel \(\rho_\tau\) bewirkt, daß \(\xi_1\) rational bekannt wird, während \(h'_\tau\) den Grad der irreduziblen Gleichung bedeutet, der \(\rho_\tau\) in dem Rationalitätsbereiche \((P,\rho_1,\ldots,\rho_\alpha,\ldots,\rho_{\tau-1},\xi_1)\) genügt. Die Multiplikation der gefundenen
\par\endgroup

\clearpage\SourceLeaf{587}\ArithmeticHead{558}{L}{TAGEBUCH MIT ERLÄUTERUNGEN.}
\par\begingroup\small
\noindent Beziehungen liefert
\[
\frac{h_\alpha\mathord{.}h_\beta\ldots h_\tau}{n}=h'_\alpha\mathord{.}h'_\beta\ldots h'_\tau.
\]
\par\noindent\hspace*{\parindent}Da sich das Produkt \(h_\alpha\mathord{.}h_\beta\ldots h_\tau\) als durch \(n\) teilbar erweist, trifft dies umsomehr für das Produkt der Grade aller Hilfsgleichungen \(X_1=0,\ X_2=0,\ldots,\ X_k=0\) zu, und damit ist unser Satz bewiesen.
\par\noindent\hspace*{\parindent}Der bequemeren Ausdrucksweise wegen wollen wir unter einer \textls[120]{Gleichungskette} \(X_1=0,\ X_2=0,\ldots,\ X_k=0\) für \(\rho_1,\rho_2,\ldots,\rho_k\) ein Gleichungssystem verstehen, bei dem \(\rho_i\) \((i=1,2,\ldots,k)\) stets Wurzel von \(X_i=0\) ist und \(X_1=0\) Koeffizienten aus dem ursprünglichen Rationalitätsbereiche \(P\), \(X_2=0\) aus dem durch \(\rho_1\) erweiterten, u.s.w. schließlich \(X_k=0\) aus dem Bereiche \((P,\rho_1,\rho_2,\ldots,\rho_{k-1})\) besitzt. Die Irreduzibilität der Gleichungen \(X_i=0\) in den Rationalitätsbereichen, denen ihre Koeffizienten angehören, wird hier \textls[120]{nicht} vorausgesetzt.
\par\noindent\hspace*{\parindent}Wir ziehen nun aus unserem Satz die folgenden Schlüsse:
\par\noindent\hspace*{\parindent}\textls[120]{Corollar I.} Die irreduzible Gleichung \(J(x)=0\) mit Koeffizienten aus \(P\) habe den Grad \(n=p_1\mathord{.}p_2\ldots p_f\), wobei \(p_1,p_2,\ldots,p_f\) lauter (auch eventuell gleiche) Primzahlen bedeuten. Es sei eine Gleichungswurzel \(\xi_1\) irgendwie als rationale Funktion von Größen \(\rho_1,\rho_2,\ldots,\rho_k\) mit Koeffizienten aus \(P\) mittels einer Gleichungskette \(Y_1=0,\ Y_2=0,\ldots,\ Y_k=0\) für \(\rho_1,\rho_2,\ldots,\rho_k\) dargestellt, und man ersetze jede Hilfsgleichung \(Y_i=0\), die in dem Rationalitätsbereiche \((P,\rho_1,\rho_2,\ldots,\rho_{i-1})\), dem ihre Koeffizienten angehören, reduzibel ist, durch die irreduzible Gleichung \(X_i=0\), der \(\rho_i\) genügt, also durch eine Gleichung von niedrigerem Grade. Dann ist für das so gebildete Gleichungssystem, das wir mit \(X_1=0,\ X_2=0,\ldots,\ X_k=0\) bezeichnen, nach unserem Satze erstens das Produkt der Grade aller Gleichungen \(X_i=0\) \((i=1,2,\ldots,k)\) durch \(p_1\mathord{.}p_2\ldots p_f\) teilbar, und zweitens enthält dieses System mindestens \(f\) Gleichungen der Grade \(p_1,p_2,\ldots,p_f\) oder wenigstens eine Gleichung, deren Grad eine zusammengesetzte Zahl ist. Diese zusammengesetzte Zahl ist wenigstens durch eine und, wenn \(k<f\) ist, sogar durch zwei der Primzahlen \(p_1,p_2,\ldots,p_f\) teilbar.
\par\noindent\hspace*{\parindent}Alles dies ergibt sich unmittelbar aus der Tatsache, daß das Produkt der Grade von \(X_1=0,\ X_2=0,\ldots,\ X_k=0\) durch \(p_1\mathord{.}p_2\ldots p_f\), den Grad von \(J(x)=0\), teilbar ist und \(p_1,p_2,\ldots,p_f\) lauter Primzahlen sind.
\par\noindent\hspace*{\parindent}\textls[120]{Corollar II.} Hat man eine Wurzel \(\xi_1\) der im Rationalitätsbereiche \(P\) irreduziblen Gleichung \(J(x)=0\) vom \(n\)-ten Grade mittels einer Gleichungskette \(X_1=0,\ X_2=0,\ldots,\ X_k=0\) für die Größen \(\rho_1,\rho_2,\ldots,\rho_k\) durch diese Größen rational mit Koeffizienten aus \(P\) dargestellt und ist das Produkt der Grade aller Gleichungen der Kette gleich \(n\), so sind die Gleichungen \(X_i=0\) \((i=1,2,\ldots,k)\) in den Bereichen, denen ihre Koeffizienten angehören, irreduzibel. Sind weiter hierbei noch die Grade aller Gleichungen der Kette Primzahlen \(p_1,p_2,\ldots,p_k\), so kann man \textls[120]{nicht unter die Grade der Gleichungen \(X_1=0,\ X_2=0,\ldots,\ X_k=0\) heruntersteigen und nicht mit einer Kette niedrigerer Gleichungen operieren.}
\par\noindent\hspace*{\parindent}Wäre nämlich, wenn das Produkt der Grade der Gleichungen der Kette gleich \(n\) ist, eine der Gleichungen \(X_i=0\) reduzibel, so könnte man sie durch einen ihrer irreduziblen Faktoren ersetzen, wodurch das Produkt der Grade der Hilfsgleichungen \(<n\) würde; dies stünde aber mit unserem Satze im Widerspruch. Mithin erweisen sich die Gleichungen \(X_i=0\) ausnahmslos als irreduzibel. Sind ferner die Grade der Gleichungen \(X_1=0,\ X_2=0,\ldots,\ X_k=0\) Primzahlen \(p_1,p_2,\ldots,p_k\) und ist \(p_1\mathord{.}p_2\ldots p_k=n\), so besagt Corollar I, daß man nicht unter die Grade \(p_1,p_2,\ldots,p_k\) der Gleichungskette \(X_1=0,\ X_2=0,\ldots,\ X_k=0\) heruntersteigen kann.
\par\noindent\hspace*{\parindent}Ist die Anzahl \(n\) der Teile, in die der Kreis geteilt werden soll, eine \textls[120]{Primzahl} und wird \(n-1\) irgendwie zerlegt in \(n-1=\alpha\mathord{.}\beta\ldots\zeta\), so führt \textsc{Gauss} im art. 352. der \textit{Disquisitiones arithmeticae}, (Werke I, S. 431) die Auflösung der Kreisteilungsgleichung \(X=0\), von der er im art. 341. (Werke I, S. 417) zeigt, daß sie irreduzibel ist, zurück auf eine Gleichungskette \(X_1=0,\ X_2=0,\ldots,\ X_k=0\) von den Graden
\par\endgroup

\clearpage\SourceLeaf{588}\ArithmeticHead{559}{R}{1801 APR. 6.}
\par\begingroup\small
\noindent\(\alpha,\beta,\ldots,\zeta\). Diese Gleichungen sind nach Corollar II in den Bereichen, denen ihre Koeffizienten angehören, irreduzibel. Wählt man die Zerlegungszahlen \(\alpha,\beta,\ldots,\zeta\) sämtlich als Primzahlen, so erfüllt die Gleichungskette von \textsc{Gauss} sämtliche Bedingungen unseres Corollars II. Mithin ist die Behauptung der \textsc{Gauss}schen Tagebuchnotiz für einen Primzahlgrad, d. h. also die Behauptung des art. 365. der \textit{Disquisitiones arithmeticae} (Werke I, S. 462) bewiesen.
\par\noindent\hspace*{\parindent}Im allgemeinen Fall, wo der Grad \(N=a^\alpha\mathord{.}b^\beta\ldots\) ist und \(a,b,\ldots\) lauter verschiedene Primzahlen bedeuten, gilt Folgendes: Die Kreisteilungsgleichung \(X=0\) vom Grade \(\varphi(N)=a^{\alpha-1}(a-1)b^{\beta-1}(b-1)\ldots\) ist irreduzibel. Ihre Auflösung wird zurückgeführt auf die Hilfsgleichungen, die bei der Teilung des Kreises in \(a,b,\ldots\) Teile auftreten und auf \(\alpha-1\) Gleichungen \(a\)-ten Grades, \(\beta-1\) Gleichungen \(b\)-ten Grades u.s.w. Man kann also die Auflösung der Gleichung \(X=0\) auf eine Gleichungskette zurückführen, die nur Gleichungen vom Primzahlgrade enthält, und für die das Produkt der Grade gleich dem Grade \(\varphi(N)\) von \(X=0\) ist. Nach unserem Corollar II ist es \textls[120]{unmöglich, die Grade der Gleichungen dieser Kette zu erniedrigen, und jede andere für die Lösung von \(X=0\) benutzbare Kette würde mindestens eine gleiche Anzahl von Gleichungen derselben Grade oder wenigstens eine Gleichung enthalten, deren Grad eine größere, zusammengesetzte Zahl ist.}
\par\noindent\hspace*{\parindent}\textsc{Gauss}' in den \textit{Disquisitiones arithmeticae} art. 366. (Werke I, S. 463) enthaltene eigene Aussage für den allgemeinen Fall (quando vero [numerus \(\varphi(N)\)] factores primos alios quam 2 puta \(p,p'\) etc. implicat, aequationes gradus \(p\)-ti, \(p'\)-ti etc. nullo modo evitari possunt) ist übrigens nicht so genau wie die obige und wie die von \textsc{Gauss} für den Fall einer Primzahlpotenz (a. a. O. tunc enim [si circulus in \(a^\alpha\) partes secandus est] praeter eas aequationes, quae ad sectionem in \(a\) partes requiruntur, necessario adhuc \(\alpha-1\) alias \(a\)-ti gradus solvere oportet; etiam has nullo modo nec evitare nec deprimere licet), wo nämlich nicht nur davon, daß die in \(\varphi(N)\) enthaltenen Primzahlen als Gradzahlen nicht vermieden bezw. nicht unterschritten werden können, sondern auch von ihrer Vielfachheit die Rede ist. Z. B. für die Teilung des Kreises in \(3^2\mathord{.}7\) Teile könnte man, da \(\varphi(N)=2^2\mathord{.}3^2\) ist, nach \textsc{Gauss} a. a. O. nur die Unvermeidlichkeit von Gleichungen zweiten und dritten Grades aussagen, während nach dem oben Bewiesenen zwei Gleichungen zweiten und zwei Gleichungen dritten Grades unvermeidlich sind. Vielleicht hängt dies damit zusammen, daß \textsc{Gauss} bei der Abfassung der \textit{Disquisitiones arithmeticae} noch nicht im Besitz des Beweises für die Irreduzibilität der \textls[120]{allgemeinen} Kreisteilungsgleichung war, den er nach Nr. 136 des \textit{Tagebuchs} erst am 12. Juni 1808 gefunden hat. Das, was \textsc{Gauss} für den allgemeinen Fall in den \textit{Disquisitiones arithmeticae} ausspricht (siehe oben), kann nämlich aus seinen Angaben für den Fall der Primzahlpotenz \textls[120]{unmittelbar} gefolgert werden. Denn, hat man die Kreisteilungsgleichung für den Fall \(N=a^\alpha\mathord{.}b^\beta\ldots\) gelöst, so hat man sie auch für die in \(N\) enthaltenen einzelnen Primzahlpotenzen \(a^\alpha,b^\beta,\ldots\) mitgelöst. Folglich sind nach \textsc{Gauss}' Angaben über die Primzahlpotenzgrade die Primfaktoren von \((a-1)a^{\alpha-1}\) bezw. \((b-1)b^{\beta-1}\) u.s.w., d. h. alle Primfaktoren von \(\varphi(N)=a^{\alpha-1}(a-1)b^{\beta-1}(b-1)\ldots\) als Gradzahlen unvermeidlich; inbezug auf die Vielfachheit versagt dieses Verfahren. Daß \textsc{Gauss} für die Primzahlpotenz auch die Vielfachheit genau bezeichnet, deutet darauf hin, daß er zur Zeit der Veröffentlichung der \textit{Disquisitiones arithmeticae} die Irreduzibilität der Kreisteilungsgleichung für einen Primzahlpotenzgrad schon gekannt haben dürfte. In der Tat läßt sich diese durch eine leichte Abänderung des im art. 341. der \textit{Disquisitiones arithmeticae} (Werke I, S. 417) enthaltenen Beweises für den Primzahlgrad erschließen, während für den allgemeinen Fall eines beliebigen Grades andere Hilfsmittel erforderlich sind.
\par\noindent\hspace*{\parindent}In den artt. 365. und 366. der \textit{Disquisitiones arithmeticae} (Werke I, S. 462, 463) hat \textsc{Gauss} aus seiner allgemeinen Aussage über die Unvermeidlichkeit der durch seine Theorie gegebenen Hilfsgleichungen den Schluß gezogen, daß der Kreis mit Zirkel und Lineal nur dann in \(N\) gleiche Teile geteilt werden kann, wenn \(\varphi(N)=2^k\) ist, d. h. wenn \(N\) von ungeraden Primteilern bloß solche von der Form \(2^m+1\) und zwar
\par\endgroup

\clearpage\SourceLeaf{589}\ArithmeticHead{560}{L}{TAGEBUCH MIT ERLÄUTERUNGEN.}
\par\begingroup\small
\renewcommand{\thefootnote}{*)}
\noindent nur in der ersten Potenz enthält. Dieses besondere Ergebnis hat zuerst \textsc{L. Wantzel}\footnote{\textsc{L. Wantzel}, \textit{Recherches sur les moyens de reconnaître si un Problème de Géométrie peut se résoudre avec la règle et le compas}, Journal de Mathém. \textbf{2}, 1837, S. 366. Der bekannte elementare Beweis, wie er von \textsc{J. Petersen}, \textit{Theorie der algebraischen Gleichungen}, Kopenhagen 1878, S. 156 und von \textsc{F. Klein}, \textit{Vorträge über ausgewählte Fragen der Elementargeometrie}, Leipzig 1895, S. 4 gegeben wird, ist eine Verschärfung des in seiner ursprünglichen Form unzureichenden Beweises von \textsc{Wantzel} auf derselben Grundlage.} zu beweisen versucht, indem er den folgenden Satz aufstellte: Soll sich eine Wurzel einer irreduziblen Gleichung durch Quadratwurzeln, d. h. durch eine Kette von Hilfsgleichungen zweiten Grades finden lassen, so muß der Gleichungsgrad notwendig eine Potenz von 2 sein. Die \textsc{Wantzel}sche Aussage folgt unmittelbar aus unserem oben bewiesenen Satze.
\par\noindent\hfill\textsc{Loewy.}
\par\endgroup
